\chapter{Unsafe Code}\label{ch-unsafe}

Zeno looks after memory in two ways: a collector, or one owner for
every value. This chapter is about the third way, the one C has: ask
for memory, use it through a \emph{pointer}, and give it back yourself.
It is for the places where firmware must control every byte, such as a
buffer that a peripheral fills, and it lives inside \verb|unsafe { }|,
a block that marks the lines Zeno cannot check.

\section{The block}

\begin{zeno}
let bytes = allocate(u8, 16)

unsafe {
    bytes[0] = 0x41
    bytes[1] = u8(bytes[0] + 1)
    print(bytes[0], bytes[1])
    free(bytes)
}
\end{zeno}

\verb|unsafe { }| is a block like the body of an \verb|if|. Inside it,
and only there, a program may do five things:

\begin{center}
\begin{tabular}{lp{7.8cm}}
\toprule
Inside unsafe only & Why Zeno cannot check it \\
\midrule
read through a pointer, \verb|p[i]| & a pointer does not know how long
  its memory is, or whether it was freed \\
write through a pointer, \verb|p[i] = v| & the same, and a wrong write
  damages something else \\
\verb|free(p)| & nothing says whether \verb|p| is still used somewhere \\
read a register, \verb|r.read()| & the address is taken on trust, and
  reading can change the hardware \\
write a register, \verb|r.write(v)| & the same \\
\bottomrule
\end{tabular}
\end{center}

Outside the block, the checker says so, and why:

\begin{zeno}[fails=yes]
let bytes = allocate(u8, 16)
print(bytes[0])
\end{zeno}

The block does not turn anything off. Types are still checked inside
it, numbers still report overflow, and lists are still looked after. It
only allows those five things. So when a program damages memory, the
lines to read are the ones inside the blocks, and \verb|grep unsafe|
finds them all.

\begin{why}
Rust showed that a language can be safe by default and still let a
driver do what drivers do, if the unsafe lines are marked and few. The
mark is a block and not a word on a function, so it can be as small as
one line. It gives no value, so \verb|let x = unsafe { p[0] }| is
refused: the block means one thing, "these lines are trusted".
\end{why}

\section{Pointers}

A pointer's type says what one item is: \verb|pointer[u8]| points at
bytes. The item is a number of one size: \verb|u8| or \verb|i8| (1
byte), \verb|u16| or \verb|i16| (2), \verb|u32| or \verb|i32| (4), or
\verb|int| or \verb|float| (8).

\verb|allocate(u8, 16)| takes the item first and then how many, and
gives a \verb|pointer[u8]| to sixteen bytes that are all zero.
\verb|p[i]| is item \verb|i|, counted from 0. \verb|p + 1| is the
address of the next \emph{item}, not the next byte, as in C.
\verb|free(p)| gives the memory back, and wants the pointer that
\verb|allocate| gave, not one worked out from it.

Everything else about a pointer is safe, and needs no block:

\begin{zeno}
let bytes = allocate(u8, 16)
let second = bytes + 1
print(second - 1 == bytes, bytes)
let all = [bytes, second]
unsafe {
    free(bytes)
}
\end{zeno}

An address that points nowhere is harmless until somebody reads or
writes through it.

\section{What zeno run catches}

A compiled program holds an address and nothing more, as a C program
does. \verb|zeno run| holds more: behind each pointer it keeps the
block, its length, and whether it was freed. So it catches the four
classic mistakes at the line that makes them, and says what a compiled
program would have done instead:

\begin{zeno}[fails=yes]
let a = allocate(u8, 16)
let b = allocate(u8, 16)
unsafe {
    a[24] = 7
    print(b[0])
}
\end{zeno}

\begin{zeno}[fails=yes]
let a = allocate(u8, 4)
unsafe {
    free(a)
    print(a[0])
}
\end{zeno}

\begin{zeno}[fails=yes]
let a = allocate(u8, 4)
unsafe {
    free(a)
    free(a)
}
\end{zeno}

Built with \verb|--no-collector|, the first program runs to its end and
prints 7: in the pool, \verb|b| lies just after \verb|a|, so item 24 of
\verb|a| is item 0 of \verb|b|. Nothing went wrong that the machine
could see. That is the whole danger, and the reason to run a program
under \verb|zeno run| first.

\section{Registers}

A register (\hyperref[ch-board]{The Board}) is made with its size and
address, \verb|register(u32, 0x4000_4400)|, and its type is
\verb|Register[u32]|. Making one and passing it around is safe;
\verb|read()| and \verb|write(v)| go inside the block. On a real board
a write is one line of C through a \verb|volatile| pointer, which tells
the C compiler that somebody else is watching this memory: it may not
skip the access, repeat it or move it.

\section{A driver}

This program from the \verb|examples| folder is a driver for a made-up
serial port. Bytes wait in a queue whose memory the program asks for
and gives back by hand, and go out through three registers. Each line
that Zeno cannot check is inside \verb|unsafe|:

\begin{zeno}[src=examples/driver.zeno]
\end{zeno}

\verb|zeno run|, the VM, and every kind of compiled build print these
same lines.

\section{Memory by hand, and the other two}

Memory from \verb|allocate| is apart from the collector and from
ownership: the collector never looks at it, and no block's end frees
it. In a build with no collector it comes from the pool, like
everything else. A list, in the same program, is still looked after
for you. Only what was asked for with \verb|allocate| is yours to free.

\begin{summary}
\verb|allocate(T, n)| gives a \verb|pointer[T]| to \verb|n| zeroed
items, and \verb|free(p)| gives it back. Reading and writing through a
pointer, \verb|free|, and reading and writing a register go inside
\verb|unsafe { }|, which allows those and changes nothing else.
Working out, comparing and keeping addresses is safe. \verb|zeno run|
catches out-of-bounds use, use after free, double frees and freeing a
moved pointer; a compiled program does not.
\end{summary}

\begin{tryit}
\begin{enumerate}
\item Allocate ten \verb|u16|s, fill them with the squares of 0 to 9,
  and print them. What happens with the square of 300?
\item Write a function that copies \verb|n| items from one pointer to
  another, and test it under \verb|zeno run| with an \verb|n| that is
  too big.
\item In \verb|driver.zeno|, find each \verb|unsafe| block and say what
  could go wrong in it.
\item Take a pointer one item on with \verb|+ 1| and \verb|free| it.
  Read the message.
\end{enumerate}
\end{tryit}
